\section{绪论}

\subsection{研究背景与意义}
脑血管病是威胁我国居民健康的重要疾病之一，其中颅内动脉瘤破裂可导致蛛网膜下腔出血，具有高致残率与高死亡率。临床上，CTA等血管成像技术已成为动脉瘤筛查与随访的重要手段，但在实际工作中，医生仍需面对海量切片、复杂血管分叉结构以及病灶边界模糊等问题，人工读片过程耗时且存在一定主观差异。

随着深度学习在医学影像分析中的快速发展，基于卷积神经网络的自动分割方法在多项器官和病灶任务中表现出显著优势。U-Net及其三维扩展结构为体数据分割提供了有效范式\cite{unet,unet3d}。同时，注意力机制与多尺度特征融合策略进一步提升了小目标与边界区域的识别能力\cite{attention_unet,nnunet}。在动脉瘤方向，已有研究开始探索从自动检测分割到破裂风险评估的一体化流程\cite{strokenet,ou2022}，但在真实工程实现中仍面临数据异构、特征融合和部署稳定性等挑战。

因此，围绕“动脉瘤精确检测与破裂风险预测”建立可复现、可扩展的工程系统，具有明确的临床辅助价值与工程实践意义。

\subsection{国内外研究现状}
\subsubsection{动脉瘤分割研究现状}
传统方法多依赖阈值分割、区域生长和形态学规则，通常需要人工调参与病例先验，泛化能力有限。深度学习方法以U-Net类网络为主流，其中3D U-Net通过体卷积直接建模三维上下文，改善了层间连续性\cite{unet3d}。Attention U-Net在跳跃连接中引入门控机制，使网络更聚焦于病灶区域\cite{attention_unet}。Kamnitsas等提出多尺度三维卷积框架以增强小病灶检出\cite{kamnitsas2017}。

面向颅内动脉瘤的研究中，Chen等采用粗到细策略提升定位精度\cite{chen2021}，Ou等验证了深度学习系统在三维血管图像中的自动诊断潜力\cite{ou2022}。然而，这些工作大多强调模型精度，对“低资源环境下如何稳定完成训练与推理”描述相对有限。

\subsubsection{破裂风险预测研究现状}
破裂风险预测通常结合临床变量、形态学指标和影像组学特征。传统统计学习方法（如逻辑回归、SVM）依赖特征工程质量，难以充分建模异构特征间的复杂关系。近年来，深度学习逐步用于风险分层任务，尤其是融合影像与表格数据的混合模型\cite{strokenet,xiong2025}。交叉注意力机制在多模态融合中表现出较强表示能力，为“数值特征+类别特征”协同建模提供了有效技术路径。

\subsection{课题来源与研究内容}
本课题来源于科研项目任务，目标是实现对颅内动脉瘤的自动化精确检测与智能化破裂风险预测。结合任务书要求与项目实现，本文完成以下工作：
\begin{enumerate}
	\item 基于三维Attention U-Net实现动脉瘤分割模型，完成数据读取、补丁采样、训练和评估流程；
	\item 构建病例级流式训练机制，通过病例按需加载与补丁化策略降低内存压力；
	\item 实现预测结果重建与NIfTI保存，形成可用于后处理和可视化的分割输出；
	\item 构建表格特征风险预测流水线，完成数据清洗、特征选择、类别编码、模型训练与阈值优化；
	\item 设计基于交叉注意力的风险预测模型，实现破裂风险二分类并输出评估报告。
\end{enumerate}

\subsection{论文结构安排}
全文共分为四章，结构如下：
\begin{itemize}
	\item 第一章为绪论，介绍研究背景、研究现状、研究内容与论文结构；
	\item 第二章阐述系统实现与方法设计，包括分割任务与风险预测任务的核心流程；
	\item 第三章给出实验设置与结果分析，说明关键配置、评价指标与结果讨论；
	\item 第四章总结全文工作并给出后续展望。
\end{itemize}
